problem:  
Let \( G \) be a compact simple Lie group and \( H \subset G \) a closed subgroup such that the homogeneous space \( M = G/H \) is connected. Suppose \( M \) admits a **\(G\)-invariant Riemannian metric with quasi-positive sectional curvature**, i.e., all sectional curvatures are nonnegative and strictly positive on a dense open subset of \(M\).  

Known classification results (Berger 1961, Wallach 1972, Aloff–Wallach 1975, Bérard-Bergery 1982) completely describe all compact homogeneous spaces admitting **strictly positive** \(G\)-invariant curvature, and subsequent work (Wilking, Petersen, Tuschmann, Grove–Ziller) has produced examples of quasi-positive curvature (some homogeneous, some not).  

**Question:**  
Is it true that for every compact simple \(G\), the existence of a \(G\)-invariant quasi-positively curved metric on \(G/H\) implies the existence of a \(G\)-invariant metric with strictly positive curvature on the same homogeneous space? In other words, does the set of \(G/H\) with quasi-positive curvature coincide with those in the known classification of \(G/H\) with positive curvature?

Equivalently, can any \(G\)-invariant quasi-positively curved metric on \(G/H\) be deformed (through \(G\)-invariant Riemannian metrics) into one with everywhere positive curvature while preserving the homogeneous structure?

Why is it a "good" problem:  
This refined problem sits precisely on the frontier of understanding the geometry of homogeneous spaces with curvature constraints. It is "good" in several dimensions:

1. **Conceptual clarity and simplicity:** The question is short and immediately comprehensible to differential geometers—does quasi-positivity within the homogeneous class actually produce new spaces or not? Yet it requires deep insight to approach.

2. **Connection to existing theory:** It directly speaks to one of the oldest classification themes—the list of positively curved homogeneous manifolds (Berger, Wallach, Aloff–Wallach, Bérard-Bergery)—and tests its rigidity. If quasi-positive curvature never enlarges this list, this would show that the classical classification is geometrically stable under relaxation of curvature conditions.

3. **Bridging analytical and algebraic structures:** Solving this would require tools from Lie theory, representation-theoretic curvature computations, and analytic deformation theory of invariant metrics (such as Cheeger deformations or Ricci flow within the invariant metric space). This bridges purely algebraic classification with geometric flow techniques.

4. **Reasonable plausibility but open status:** Known constructions (by Wilking and collaborators) yield quasi-positive curvature via \(G\)-invariant Cheeger deformations, but no known examples establish quasi-positive curvature on a homogeneous space failing to support positive curvature. Thus, the conjecture is plausibly open and fundamental.

5. **Potential outcome:** A positive resolution would imply deep rigidity—homogeneous quasi-positive curvature cannot occur “on new spaces.” A negative answer would produce previously unknown homogeneous spaces with rich geometric properties, redefining the boundary of the classification program.

Hence, the question is precise, bridges established theories, and, though simple to state, would demand substantial conceptual innovation to resolve—hallmarks of a “good” research problem.